%%%%%%%%%%%%%%%%%%%%%%%%%%%%%%%%%%%%%%%%%%%%%%%%%%%%%%%%%%%%%%%%%%%%%%
% Lớp 12 | Ứng dụng đạo hàm để khảo sát hàm số | Đường tiệm cận
% Dạng 2D1?4-2: Bài toán xác định các đường tiệm cận của hàm số có chứa tham số
% Số câu: 33  (sinh tự động bởi tach_cau.py; từ nay thầy sửa trực tiếp file này)
%%%%%%%%%%%%%%%%%%%%%%%%%%%%%%%%%%%%%%%%%%%%%%%%%%%%%%%%%%%%%%%%%%%%%%

%%%%%%%%%%%%%%%% TRẮC NGHIỆM 4 PHƯƠNG ÁN %%%%%%%%%%%%%%%%

% nguồn: Phan1_Lop_12_Dai_so_Chuong_1
\begin{ex}%[2D1H4-2]%[ma:45b7fc8c]
	Số giá trị  nguyên của tham số $m$ sao cho đồ thị hàm số $y=\dfrac{\sqrt{x+1}}{x^2-2x+m}$ có đúng ba đường tiệm cận.
	 \choice
		 {$5$}
		{Vô  số}
		 {\True $3$}
		 {$4$}
	\loigiai{ 
		$\lim\limits_{x\to+\infty} y=\lim\limits_{x\to+\infty}\dfrac{\sqrt{x+1}}{x^2-2x+m}=0$ nên đồ thị hàm số có tiệm  cận ngang $y=0$.\\
		Vậy để đồ thị hàm số $y=\dfrac{\sqrt{x+1}}{x^2-2x+m}$ có đúng ba  đường tiệm cận thì đồ thị hàm số có hai tiệm cận đứng $\Leftrightarrow x^2-2x+m=0$ có 2 nghiệm phân biệt $x_1, x_2$ lớn hơn hoặc bằng $-1$ \\
	$ \Leftrightarrow\heva{&\Delta’>0\\&(x_1+1)(x_2+1)\geq  0\\&(x_1+1)+(x_2+1)>0}\Leftrightarrow\heva{&1-m>0\\&m+2+1>0\\&4>0}\Leftrightarrow-3<m<1 $. Vì $m\in\mathbb{Z}\Rightarrow m\in\{-2;-1;0\}$.}
\end{ex}

% nguồn: Phan1_Lop_12_Dai_so_Chuong_1
\begin{ex}%[2D1H4-2]%[ma:2675f2ce]
	Tìm tất cả các giá trị của tham  số $m$ để đồ thị hàm số $y=\dfrac{2x+4}{x-m}$ có tiệm cận đứng?
	\choice 
		{$m >-2$} 
		 {$m=-2$}
		 {$m <-2$}
		{\True  $m\neq-2$}
	 \loigiai{
		Để $x=m$ là tiệm cận đứng của đồ thị hàm số  $y=\dfrac{u(x)}{v(x)}=\dfrac{2x+4}{x-m}$ thì $$\heva{&v(m)=0\\&u(m)\neq 0}
	\Leftrightarrow\heva{&m-m=0\\&2m+4\neq 0}\Leftrightarrow\heva{&0=0\\&m\neq-2}\Leftrightarrow m\neq-2.$$} 
\end{ex}

% nguồn: Phan1_Lop_12_Dai_so_Chuong_1
\begin{ex}%[2D1H4-2]%[ma:97dd0a0f]
	Giá trị $m$ để tiệm cận đứng của  đồ thị hàm số $y=\dfrac{2x+2m-1}{x+m}$ đi qua điểm $M(3;1)$ là
	 \choice
		{\True  $m=-3$}
		{$m=-1$} 
		 {$m=2$}
		 {$m=3$}
	 \loigiai{
		Tiệm cận đứng của đồ thị hàm số đi qua điểm $M(3;1)$ nên đồ thị hàm có tiệm  cận đứng là $x=3$.\\
		Suy ra $x+m=0$ có nghiệm là $3$ do vậy $3+m=0\Leftrightarrow m=-3$.\\ 
		Thử lại, với $m=-3\Rightarrow y=\dfrac{2x-7}{x-3}$ có $\lim\limits_{x\to 3^+}  y=\lim\limits_{x\to 3^+}\dfrac{2x-7}{x-3}=-\infty$ và $\lim\limits_{x\to 3^-} y=\lim\limits_{x\to 3^-}\dfrac{2x-7}{x-3}=+\infty$.\\
	Vậy  $m=-3$.}
\end{ex}

% nguồn: Phan1_Lop_12_Dai_so_Chuong_1
\begin{ex}%[2D1H4-2]%[ma:7f6bcd1f]
	Cho đường cong $(C)\colon y=\dfrac{x+a}{x-1}$. Biết điểm $M$ thuộc $(C)$. Tính tổng các giá trị của tham số $a$ để tiếp tuyến của  $(C)$ tại $M$ tạo với hai đường tiệm cận của $(C)$ một tam giác có diện tích bằng $3\sqrt{2}+2$.
	 \choice
		{$3\sqrt{2}-2$} 
		 {$6$}
		{\True  $-2$}
		{$\dfrac{3\sqrt{2}-4}{2}$} 
	\loigiai{ 
		Tiệm cận đứng $x=1$; tiệm cận ngang $y=1$. Ta  có $f’(x)=\dfrac{-1-a}{(x-1)^2}$.\\
		Phương trình tiếp  tuyến tại $M$ của $(C)$ là $(d)\colon y=\dfrac{-1-a}{(x_0-1)^2}(x-x_0)+\dfrac{x_0+a}{x_0-1}$.\\
		Giao điểm của $(d)$ với tiệm cận  đứng: $A\left(1;\dfrac{x_0+2a+1}{x_0-1}\right)$.\\
		Giao  điểm của $(d)$ với tiệm cận ngang: $B(2x_0-1;1)$.\\
		 Giao hai đường tiệm cận của $(C)$: $I(1;1)$. Khi đó $IA=\left|\dfrac{2a+2}{x_0-1}\right|$; $IB=2|x_0-1|$.\\
		 Ta có $S_{IAB}=\dfrac{1}{2}IA\cdot IB =|2a+2|=3\sqrt{2}+2\Leftrightarrow\hoac{&a=\dfrac{3\sqrt{2}}{2}\\&a=\dfrac{-3\sqrt{2}-4}{2}.}$ \\
	 Vậy tổng cần tìm là $-2$.}
\end{ex}

% nguồn: Phan1_Lop_12_Dai_so_Chuong_1
\begin{ex}%[2D1H4-2]%[ma:39a0bdce]
	Cho hàm số $y=\dfrac{x-2}{x^2-2mx-m-2}$. Biết với $m=\dfrac{a}{b}$ ($a,b\in\mathbb{N}$,  $\dfrac{a}{b}$ tối giản) thì đồ thị hàm số có đúng $2$ đường tiệm cận. Tính $a+b$.
	 \choice
		{\True $a+b=7$} 
		{\True $a+b=5$} 
		{$a+b=8$} 
		 {$a+b=6$}
	 \loigiai{
		Để đồ thị hàm số có đúng $2$ đường tiệm cận thì hoặc phương trình $x^2-2mx-m-2=0$ có nghiệm kép $x=2$ hoặc phương trình  $x^2-2mx-m-2=0$ phải có hai nghiệm (một nghiệm $x_1=2$ và một nghiệm $x_2\neq 2$).\\
		Do $\Delta’=m^2+m+2>0,\forall m$ nên ta chỉ xét trường hợp thứ hai phương trình $x^2-2mx-m-2=0$ có hai nghiệm phân biệt. Thay $x=2$ vào phương trình ta được  $m=\dfrac{2}{5}$ (thỏa mãn).\\
	Vậy $a=2,b=5,a+b=7$.} 
\end{ex}

% nguồn: Phan1_Lop_12_Dai_so_Chuong_1
\begin{ex}%[2D1H4-2]%[ma:7ebb86af]
	Tìm tất cả các giá  trị của tham số $m$ sao cho đồ thị hàm số $y=\dfrac{x-1}{x^3+3x^2+m+1}$ có đúng một tiệm cận đứng.
	 \choice
		 {$\hoac{&m\leq-4\\&m>0}$}
		 {$\hoac{&m <-5\\&m >-1}$}
		 {$-5\leq m <-1$}
		 {\True $\hoac{&m\leq-5\\&m >-1}$}
	 \loigiai{
		Đặt $f(x)=x^3+3x^2+m+1$;  $f’(x)=3x^2+6x; f’(x)=0\Leftrightarrow\hoac{&x=-2\\&x=0.}$ \\
		\textbf{Trường hợp 1:} $g(1)=0\Rightarrow 1+3+m+1=0\Leftrightarrow  m=-5$.\\
		Khi đó $g(x)=x^3+3x^2-4=(x-1)(x+2)^2$  \\
		$ \Rightarrow y=\dfrac{x-1}{(x-1)(x+2)^2}\Rightarrow\lim\limits_{x\to  1} y=\lim\limits_{x\to 1}\dfrac{x-1}{(x-1)(x+2)^2}=\dfrac{1}{9} $;\\
		$\Rightarrow\lim\limits_{x\to-2^+} y=\lim\limits_{x\to-2^+}\dfrac{x-1}{(x-1)(x+2)^2}=+\infty$;  $\lim\limits_{x\to-2^-} y=\lim\limits_{x\to-2^-}\dfrac{x-1}{(x-1)(x+2)^2}=+\infty$.\\
		Nên đồ thị hàm số chỉ có một tiệm cận  đứng $x=-2$.\\
		 \textbf{Trường hợp 2:} $g(1)\neq 0\Rightarrow 1+3+m+1\neq 0\Leftrightarrow m\neq-5$.\\
		Đồ thị hàm số đã cho có đúng một tiệm cận đứng khi phương trình $g(x)=0$ có nghiệm duy nhất $\Leftrightarrow g(0)\cdot g(-2)>0\Leftrightarrow(m+1)(m+5)>0\Leftrightarrow\hoac{&m <-5\\&m  >-1.}$ \\
	Kết hợp hai trường hợp ta có  những giá trị $m$ cần tìm là $\hoac{&m\leq-5\\&m >-1}$.}
\end{ex}

% nguồn: Phan1_Lop_12_Dai_so_Chuong_1
\begin{ex}%[2D1H4-2]%[ma:a80a8ac7]
	Có bao nhiêu giá  trị của tham số để đồ thị hàm số $y=\dfrac{mx^2-1}{x^2-3x+2}$ có đúng $2$ tiệm cận?
	\choice 
		{\True  $2$}
		 {$1$}
		{$4$} 
		{$3$} 
	 \loigiai{
		Ta có  $\lim\limits_{x\to-\infty} y=\lim\limits_{x\to-\infty}\dfrac{mx^2-1}{x^2-3x+2}=\lim\limits_{x\to-\infty}\dfrac{m-\dfrac{1}{x^2}}{1-\dfrac{3}{x}+\dfrac{2}{x^2}}=m\Rightarrow y=m$ là tiệm cận ngang của đồ thị hàm số.\\
		Để đồ thị hàm số $y=\dfrac{mx^2-1}{x^2-3x+2}$ có đúng $2$ tiệm cận $\Leftrightarrow y=\dfrac{mx^2-1}{x^2-3x+2}$ có đúng $1$ tiệm cận đứng. Ta có $x^2-3x+2=0\Leftrightarrow\hoac{&x=1\\&x=2}\Rightarrow\hoac{&m{\cdot 1}^2-1=0\\&m{\cdot 2}^2-1=0}\Rightarrow\hoac{&m=1\\&m=\dfrac{1}{4}.}$  \\
	Vậy có $2$ giá trị tham số thỏa mãn  yêu cầu đề bài.}
\end{ex}

% nguồn: Phan1_Lop_12_Dai_so_Chuong_1
\begin{ex}%[2D1H4-2]%[ma:1d1a5203]
	Số các giá trị nguyên của tham số $m$  thuộc $[-2025; 2025]$ để đồ thị hàm số $y=\dfrac{2x+4}{x-m}$ có tiệm cận đứng nằm bên trái trục tung là
	 \choice
		{$2020$} 
		{\True $2025$} 
		 {$4041$}
		 {$4042$}
	 \loigiai{
		Đồ thị hàm số $y=\dfrac{2x+4}{x-m}$ có tiệm cận  đứng là đường thẳng $x=m$ nằm bên trái trục tung $\Leftrightarrow m<0$.\\
		Mặt  khác $\heva{&m\in[-2025; 2025]\\&m\in\mathbb{Z}}\Rightarrow m\in\left\{-2025;-2020;\ldots;-1\right\}$.\\
	Vậy có $2025$ giá trị nguyên của $m$ thoả mãn yêu cầu.} 
\end{ex}

% nguồn: Phan1_Lop_12_Dai_so_Chuong_1
\begin{ex}%[2D1H4-2]%[ma:c923bffa]
	Cho đồ thị hàm số $y=\dfrac{(2m-n)x^2+mx+1}{x^2+mx+n-6}$  nhận trục hoành và trục tung làm hai tiệm cận. Giá trị $m+n$ là
	 \choice
		{$8$} 
		{\True $9$} 
		 {$6$}
		{$-6$} 
	 \loigiai{
		 Điều kiện: $x^2+mx+n-6\neq 0$.\\
		Phương trình đường tiệm cận ngang của  đồ thị hàm số là $y=2m-n$.\\
		Vì đồ thị hàm số nhận trục hoành làm tiệm cận ngang nên $2m-n=0$.\\ 
		 Đặt $f(x)=(2m-n)x^2+mx+1$ và $g(x)=x^2+mx+n-6$.\\
	Vì $f(0)\neq 0$ với mọi $m$, $n$ nên đồ thị nhận trục tung $x=0$ là tiệm cận đứng khi $g(0)=0\Leftrightarrow n=6$.  Suy ra $m=\dfrac{n}{2}=3$. Vậy $m+n=9$.}
\end{ex}

% nguồn: Phan1_Lop_12_Dai_so_Chuong_1
\begin{ex}%[2D1H4-2]%[ma:b6926885]
	Có bao nhiêu giá trị nguyên $m\in[-10;10]$ sao cho đồ thị hàm số $y=\dfrac{x-1}{x^2+4x-m-3}$ có  hai đường tiệm cận đứng?
	\choice 
		 {$19$}
		 {$15$}
		{\True  $16$}
		 {$17$}
	 \loigiai{
		Ta có đồ thị hàm  số $y=\dfrac{x-1}{x^2+4x-m-3}$ có hai đường tiệm cận đứng khi phương trình $x^2+4x-m-3$ có hai nghiệm phân biệt khác $1\Leftrightarrow\heva{&2^2-(-m-3)>0\\&1^2+4\cdot 1-m-3\neq 0}\Leftrightarrow\heva{&m >-7\\&m\neq 2.}$ \\
		 Từ đó ta suy ra $m=\left\{-6,-5,-4,-3,-2,-1,0,1,3,4,5,6,7,8,9,10\right\}$.\\
	Vậy có $16$ giá trị nguyên của $m$ thỏa  mãn.}
\end{ex}

% nguồn: Phan1_Lop_12_Dai_so_Chuong_1
\begin{ex}%[2D1H4-2]%[ma:5e754b93]
	Cho hàm  số $y=\dfrac{x-2}{mx^2-2(m-1)x+m-2}$ có đồ thị $(C_m)$. Có bao nhiêu giá trị nguyên của tham số $m$ để $(C_m)$ có đúng $2$ đường tiệm cận?
	 \choice
		{\True $2$} 
		{$1$} 
		{$0$} 
		{$3$} 
	\loigiai{ 
		\textbf{Trường hợp 1:} $m=0$, ta được hàm số $y=\dfrac{x-2}{2x-2}$, suy ra đồ thị hàm số có tiệm cận  đứng $x=1$; tiệm cận ngang $y=\dfrac{1}{2}$ (thỏa mãn).\\
		\textbf{Trường hợp 2:} $m\neq 0$, ta có  $\lim\limits_{x\to+\infty} y=\lim\limits_{x\to-\infty} y=0\Rightarrow y=0$ là tiệm cận ngang của đồ thị hàm số.\\
		Xét phương  trình: $mx^2-2(m-1)x+m-2=0(m\neq 0)\Leftrightarrow\hoac{&x=1\\&x=\dfrac{m-2}{m}.}$ \\
		Để đồ thị hàm số $(C_m)$ có đúng $2$ đường tiệm cận thì đồ thị hàm  số $(C_m)$ chỉ có $1$ tiệm cận đứng\\
		$ \Leftrightarrow mx^2-2(m-1)x+m-2=0 $ hoặc có hai nghiệm phân biệt trong đó có một nghiệm bằng 2 hoặc có  nghiệm kép $\Leftrightarrow\hoac{&\dfrac{m-2}{m}=2\\&\dfrac{m-2}{m}=1}\Leftrightarrow m=-2$.\\
	Vậy có $2$ giá trị nguyên của tham số $m$ thỏa yêu  cầu bài toán.}
\end{ex}

% nguồn: Phan1_Lop_12_Dai_so_Chuong_1
\begin{ex}%[2D1V4-2]%[ma:5ae884bb]
	Tìm tổng tất cả các giá trị của tham số thực $m$ để đồ thị  hàm số $y=\dfrac{x-1}{x-m}$ có hai đường tiệm cận tạo với hai trục tọa độ một hình chữ nhật có diện tích bằng $5$.
	 \choice
		 {$2$}
		{$4$} 
		{\True $0$} 
		 {$5$}
	\loigiai{ 
		Xét hàm nhất biến $y=\dfrac{x-1}{x-m}$ có tiệm cận đứng  $x=m$ và tiệm cận ngang $y=1$.\\
		Để hai đường tiệm cận tạo với hai  trục tọa độ một hình chữ nhật có diện tích bằng $5$.\\
		khi và chỉ khi: $|m|\cdot  1=5\Leftrightarrow\hoac{&m=5\\&m=-5.}$ \\
		Vậy có hai giá  trị $m$ thỏa mãn và tổng chúng bằng $0$.
	 }
\end{ex}

% nguồn: Phan1_Lop_12_Dai_so_Chuong_1
\begin{ex}%[2D1V4-2]%[ma:8a618f1b]
	Cho hàm số $y=\dfrac{x^2+(m-1)x+m^2-2m+1}{1-x}\quad (1)$, $m$ là tham số  thực. Tìm giá trị của $m$ để đường tiệm cận xiên của đồ thị hàm số $(1)$ tạo với hai trục tọa độ một tam giác có diện tích bằng $\dfrac{1}{2}$.
	\choice 
		{$m=\pm  2$}
		{$m=\pm  3$}
		{$m=\pm 4$} 
		{\True $m=\pm  1$}
	\loigiai{ 
		Tập xác định: $\mathscr{D}=\mathbb{R}\setminus\{1\}$. Ta có  $y=\dfrac{x^2+(m-1)x+m^2-2m+1}{1-x}=-x-m+\dfrac{m^2-m+1}{1-x}$.\\
		Đồ thị hàm số có tiệm cận xiên là  $\Delta\colon y=-x-m$.\\
		Tiệm cận xiên cắt hai trục tọa độ tại hai điểm $A(0;-m)$ và  $B(-m;0)$.\\
		Diện  tích tam giác $OAB$ là $S=\dfrac{1}{2}OA\cdot OB=\dfrac{1}{2}|y_A|\cdot|y_B|=\dfrac{1}{2}m^2$.\\
		Theo giả thiết $S=\dfrac{1}{2}\Leftrightarrow  m^2=1\Leftrightarrow m=\pm 1$.\\
	 Vậy giá trị $m$ cần tìm là $m=\pm 1$.}
\end{ex}

% nguồn: Phan1_Lop_12_Dai_so_Chuong_1
\begin{ex}%[2D1C4-2]%[ma:b250b8b3]
	 Tìm tổng tất cả các giá trị nguyên dương của tham số $m$ sao cho đồ thị hàm số $y = \dfrac{x+2}{x^2-4x+m}$ có 3 đường tiệm cận?
	 \choice
		{\True $6$} 
		 {$2$}
		 {$5$}
		 {$4$}
	\loigiai{Ta có $\lim\limits_{x \to \pm \infty} y  = 0$. Suy ra $y=0$ là phương trình tiệm cận ngang của đồ thị hàm số.\\
		Để hàm số có 3 đường  tiệm cận thì phương trình $x^2-4x+m=0$ có 2 nghiệm phân biệt với mọi $x \ne -2$.\\
		Suy ra $\heva{&\Delta = 16 - 4m > 0\\&12 + m \ne 0} \Leftrightarrow m  \in (-\infty; 4) \setminus \{-12\}$.\\
		Vậy tổng  các giá trị nguyên dương của $m$ thỏa $3+2+1=6$.
	} 
\end{ex}


%%%%%%%%%%%%%%%% ĐÚNG/SAI %%%%%%%%%%%%%%%%

% nguồn: Phan1_Lop_12_Dai_so_Chuong_1
\begin{ex}%[2D1H4-2]%[ma:b7d8afa4]
	Cho hàm số $y=\dfrac{mx^2+6x-2}{x+2}$, với $m$ là tham số. 
	 \choiceTF
		{\True Tập xác định của  hàm số là $\mathscr{D}=\mathbb{R}\setminus\{-2\}$}
		 {Đồ thị hàm số có tiệm cận ngang khi $m>0$}
		{Đồ thị hàm số có tiệm cận đứng khi  $m\neq 0$}
		{\True Tập hợp tất cả các giá trị của $m$ để đồ thị hàm số có hai đường tiệm cận là  $\mathbb{R}\setminus\left\{\dfrac{7}{2}\right\}$}
	\loigiai{ 
		\begin{itemchoice}
			 \itemch Điều kiện $x+2\neq 0\Leftrightarrow x\neq-2$.\\
			Vậy tập xác định  là $\mathscr{D}=\mathbb{R}\setminus\{-2\}$.
			 \itemch Đồ thị hàm số có tiệm cận ngang khi hệ số $x^2$ trên tử phải bằng $0$ suy ra $m=0$.
			\itemch Đồ thị hàm số có tiệm cận  đứng khi $x=-2$ không là nghiệm của tam thức.\\
			$g(x)=mx^2+6x-2\Leftrightarrow g(-2)\neq  0\Leftrightarrow m\neq\dfrac{7}{2}$.
			\itemch Đồ thị hàm số có chắc chắn một tiệm cận xiên (hoặc ngang).\\ 
			Suy  ra để đồ thị hàm số có hai đường tiệm cận thì nó phải có một tiệm cận đứng.\\
			Điều  này tương đương với $m\neq\dfrac{7}{2}$.
	\end{itemchoice}}
\end{ex}

% nguồn: Phan1_Lop_12_Dai_so_Chuong_1
\begin{ex}%[2D1H4-2]%[ma:5e43b9fa]
	Gọi  $(C)$ là đồ thị của hàm số $y=\dfrac{(m+1)x^2+(2m+1)x+m+2}{x+1}$.
	\choiceTF 
		{Tiệm cận đứng của đồ thị hàm số  là $x=1$}
		{\True Tiệm cận xiên của đồ thị hàm số đi qua điểm  $M(1;2)$ khi $m=\dfrac{1}{2}$}
		{Tiệm cận xiên của đồ thị hàm  số vuông góc với đường thẳng $d\colon 3x+4y-5=0$ khi $m=1$}
		{\True Giao điểm của hai đường tiệm cận của $(C)$  luôn thuộc Parabol $(P)\colon y=-x^2$}
	\loigiai{ 
		\begin{itemchoice}
			\itemch Tiệm cận đứng của đồ thị  hàm số là $x=-1$.
			 \itemch $y=\dfrac{(m+1)x^2+(2m+1)x+m+2}{x+1}=(m+1)x+m+\dfrac{2}{x+1}$.\\
			Đồ thị hàm số có tiệm cận xiên là $\Delta\colon  y=(m+1)x+m$.\\
			Vì tiệm cận xiên đi qua điểm $M(1;2)$ nên $2=(m+1)\cdot 1+m\Leftrightarrow 2m=1\Leftrightarrow  m=\dfrac{1}{2}$.
			\itemch Ta có $\Delta\colon y=(m+1)x+m\Leftrightarrow(m+1)x-y+m=0$.\\ 
			Tiệm cận xiên vuông góc với đường  thẳng $d\colon 3x+4y-5=0$.\\
			nên $(m+1)\cdot 3+4\cdot  (-1)=0\Leftrightarrow m+1=\dfrac{4}{3}\Leftrightarrow m=\dfrac{1}{3}$.
			\itemch  Giao điểm của hai đường tiệm cận là $I(-1;-1)\in(P)$.
	\end{itemchoice}}
\end{ex}

% nguồn: Phan1_Lop_12_Dai_so_Chuong_1
\begin{ex}%[2D1H4-2]%[ma:d38d7da5]
	Cho hàm số $y=f(x)=\dfrac{nx^2+1}{mx+1} (mn\neq 0)$  có đồ thị $(C)$ như hình vẽ:
	\begin{center}
		\begin{tikzpicture}[scale=1,font=\footnotesize,line join=round,line cap=round,>=stealth]
			\def\a{1} \def\b{0} \def\c{1} \def\d{2} \def\e{1} 
			\def\xmin{-6}  \def\xmax{5}
			\def\ymin{-5}  \def\ymax{5}
			\draw[->] (\xmin,0)--(\xmax,0) node [below]{$x$}; 
			 \draw[->] (0,\ymin)--(0,\ymax) node [right]{$y$};
			 \node at (0,0) [above right]{$O$};
			\clip (\xmin+0.1,\ymin+0.1) rectangle (\xmax-0.1,\ymax-0.1); 
			 \draw[smooth,samples=300,domain=\xmin:(-\e/\d-0.1)] plot (\x,{(\a*(\x)^2+\b*(\x)+\c)/(\d*(\x)+\e)});
			\draw[smooth,samples=300,domain=(-\e/\d+0.1:\xmax)] plot (\x,{(\a*(\x)^2+\b*(\x)+\c)/(\d*(\x)+\e)}); 
			\draw[smooth,samples=300,domain=\xmin:\xmax] plot (\x,{(0.5*(\x)-0.25)}); 
			\draw  (-\e/\d,\ymin)--(-\e/\d,\ymax);
			\fill (-0.5,0) circle (1.0pt) node[above  left]{$-\dfrac{1}{2}$};
			\node at  (-4.1,-1)[below]{$y=\dfrac{1}{2}x-\dfrac{1}{4}$};
		\end{tikzpicture}
	\end{center}
	\choiceTF 
		 {$\dfrac{n}{m}=-\dfrac{1}{4}$}
		 {$m=-\dfrac{1}{2}$}
		 {\True $m+n=3$}
		{$\lim\limits_{x\to+\infty}\left[f(x)-\dfrac{1}{2}x\right]=0$} 
	\loigiai{ 
		\begin{itemchoice}
			\itemch $\lim\limits_{x\to+\infty}\dfrac{f(x)}{x}=\lim\limits_{x\to+\infty}\dfrac{nx^2+1}{mx^2+x}=\dfrac{n}{m}=\dfrac{1}{2}$. 
			\itemch Đường tiệm cận đứng của đồ thị là $x=-\dfrac{1}{m}=-\dfrac{1}{2}\Rightarrow m=2$. 
			\itemch Từ hai ý a) và b) ta có $\heva{&m=2n\\&m=2}\Rightarrow\heva{&m=2\\&n=1}\Rightarrow  m+n=3$; $n=2\Rightarrow m+n=3$.\\
			\itemch Đường  tiệm cận xiên $y=\dfrac{1}{2}x-\dfrac{1}{4}\Rightarrow\lim\limits_{x\to+\infty}\left[f(x)-\dfrac{1}{2}x\right]=-\dfrac{1}{4}$.
		\end{itemchoice}
	 }
\end{ex}

% nguồn: Phan1_Lop_12_Dai_so_Chuong_1
\begin{ex}%[2D1V4-2]%[ma:12d90ce6]
	Cho hàm số $y=\dfrac{mx^2+\left(3m^2-2\right)x-2}{x+3m}$ có  đồ thị là $(C)$ với $m$ là tham số.
	\choiceTF 
		 {\True Đồ thị hàm số có hai đường tiệm cận khi $m\neq\dfrac{1}{3}$}
		 {\True Khi $m=1$ thì đường tiệm cận xiên của đồ thị hàm số có phương trình $y=x-2$}
		{Biết đồ thị $(C)$ có hai đường  tiệm cận và khi $m=\pm 2$ thì góc tạo bởi hai đường tiệm cận của đồ thị hàm số bằng $45^{\circ}$}
		{\True Gọi $S$ là tập các giá trị của tham số $m$ để đồ thị $(C)$ có đường tiệm cận xiên cắt các trục toạ độ tại hai điểm $A$ và $B$ sao cho tam giác $OAB$ có diện tích bằng $4$.  Khi đó tổng bình phương các phần tử của tập $S$ bằng $4$}
	 \loigiai{
		Ta có $y=mx-2+\dfrac{6m-2}{x+3m}$. 
		\begin{itemchoice}
			\itemch Đồ  thị hàm số có hai đường tiệm cận khi và chỉ khi $6m-2\neq 0\Leftrightarrow m\neq\dfrac{1}{3}$.
			\itemch Khi $m=1$ thì  đường tiệm cận xiên của đồ thị hàm số có phương trình $y=x-2$.\\
			Phương trình hai  đường tiệm cận là\\
			$\Delta_1\colon x=-3m\Leftrightarrow  x+3m=0;\quad\Delta_2\colon y=mx-2\Leftrightarrow mx-y-2=0$.
			\itemch Véc-tơ pháp tuyến của $\Delta_1$ và $\Delta_2$ lần lượt là $\overrightarrow{n}_1=(1; 0)$; $\overrightarrow{n}_2=(m;-1)$.\\ 
			Góc tạo bởi hai đường tiệm cận của đồ thị hàm số bằng $45^{\circ}$ khi và chỉ  khi $$ \cos45^{\circ}=\dfrac{\left|\overrightarrow{n}_1\cdot\overrightarrow{n}_2\right|}{\left|\overrightarrow{n}_1\right|\left|\overrightarrow{n}_2\right|}
			\Leftrightarrow\dfrac{|m|}{\sqrt{m^2+1}}=\dfrac{\sqrt{2}}{2}\Leftrightarrow  2m^2=m^2+1\Leftrightarrow m=\pm 1.$$
			\itemch Hàm số có tiệm cận xiên khi $m\neq 0$ và $m\neq\dfrac{1}{3}$. Khi  đó $A(0;-2); B\left(\dfrac{2}{m}; 0\right)$.\\
			Ta có diện tích tam giác $OAB=\dfrac{1}{2}\cdot OA\cdot OB=4\Leftrightarrow\dfrac{1}{2}\cdot|-2|\cdot\left|\dfrac{2}{m}\right|=4\Leftrightarrow m=\pm 2$ nên tổng bình  phương các phần tử của tập $S$ bằng $4$.
	\end{itemchoice}}
\end{ex}

% nguồn: Phan1_Lop_12_Dai_so_Chuong_1
\begin{ex}%[2D1V4-2]%[ma:907157a5]
	Cho  hàm số $y=\dfrac{mx^2+\left( 3m^2-2 \right)x-2}{x+3m}$ có đồ thị là $\left( C \right)$ với $m$ là tham số.
	\choiceTF 
		{\True Khi $m=1$, đồ thị hàm số  có đường tiệm cận đứng là $x=-3$}
		{\True Khi $m=1$ thì đường tiệm cận xiên của đồ thị hàm số có phương  trình $y=x-2$}
		{Đồ thị hàm  số có hai đường tiệm cận khi $m\ne -3$}
		{Gọi $S$ là tập các giá trị của tham số $m$ để đồ thị $\left( C \right)$ có đường tiệm cận xiên cắt các trục toạ độ tại hai điểm $A$ và $B$ sao cho tam giác $OAB$ có diện tích  bằng $4$. Khi đó tổng bình phương các phần tử của tập $S$ bằng $4$}
	\loigiai{ 
		\begin{itemchoice}
			\itemch {\bf Đúng.}\\ 
			Khi $m=1$, đồ thị hàm số $y=\dfrac{x^2+x-2}{x+3}$ có đường tiệm cận đứng  là $x=-3$.
			\itemch {\bf Đúng.}\\ 
			Khi $m=1$, $y=x-2+\dfrac{4}{x+3}$ thì đường tiệm cận xiên của  đồ thị hàm số là $y=x-2$.
			\itemch  {\bf Sai.}\\
			Ta  có: $y=mx-2+\dfrac{6m-2}{x+3m}$.\\
			Đồ thị hàm số  có hai đường tiệm cận khi và chỉ khi $6m-2\ne 0\Leftrightarrow m\ne \dfrac{1}{3}$.\\
			Khi đó phương  trình hai đường tiệm cận là: $\Delta_1\colon x=-3m$ và $\Delta_2 \colon y=mx-2$.
			 \itemch {\bf Sai.}\\
			Hàm số có tiệm cận xiên khi $m\ne 0$ và $m\ne \dfrac{1}{3}$. Khi đó $A\left( 0; -2 \right); B\left( \dfrac{2}{m}; 0 \right)$.\\ 
			Ta có diện tích tam giác $OAB=\dfrac{1}{2}\cdot OA\cdot  OB=4\Leftrightarrow \dfrac{1}{2}\cdot \left| -2 \right|\cdot \left| \dfrac{2}{m} \right|=4\Leftrightarrow m=\pm \dfrac{1}{2}$ nên tổng bình phương các phần tử của tập $S$ bằng $\dfrac{1}{2}$.
		\end{itemchoice}
	 }
\end{ex}


%%%%%%%%%%%%%%%% TRẢ LỜI NGẮN %%%%%%%%%%%%%%%%

% nguồn: Phan2_Lop_12_Dai_so_Chuong_1
\begin{bt}%[2D1H4-2]%[ma:fc8fbcf7]
	Có bao  nhiêu giá trị của $m$ để đồ thị hàm số $y=\dfrac{mx^2-1}{x^2-3x+2}$ có đúng $2$ đường tiệm cận?
	 \shortans{2}
	\loigiai{ 
		Ta có $\lim\limits_{x\to\pm\infty} y=\lim\limits_{x\to\pm\infty}\dfrac{m-\dfrac{1}{x^2}}{1-\dfrac{3}{x}+\dfrac{2}{x^2}}=m\Rightarrow$ tiệm cận ngang  $y=m$.\\
		Để hàm số có đúng $2$ đường tiệm cận thì hàm số  có đúng $1$ tiệm cận đứng.\\
		 Suy ra $mx^2-1=0$ có một nghiệm bằng $1$ hoặc bằng $2$.\\
		 Khi đó $\hoac{&m-1=0\\&4m-1=0}\Leftrightarrow\hoac{&m=1\\&m=\dfrac{1}{4}.}$\\
		 Với $m=1\Rightarrow y=\dfrac{x^2-1}{x^2-3x+2}=\dfrac{x+1}{x-2}\Rightarrow\lim\limits_{x\to 2^+} y=+\infty\Rightarrow$ tiệm cận đứng $x=2$.\\
		Với $m=\dfrac{1}{4}\Rightarrow y=\dfrac{\dfrac{1}{4}x^2-1}{x^2-3x+2}=\dfrac{x+2}{4(x-1)}\Rightarrow\lim\limits_{x\to 1^+} y=+\infty\Rightarrow$ tiệm  cận đứng $x=1$.\\
	Vậy có $2$ giá trị $m$ thỏa bài.} 
\end{bt}

% nguồn: Phan2_Lop_12_Dai_so_Chuong_1
\begin{bt}%[2D1H4-2]%[ma:4a0aa5bd]
	Biết đồ thị hàm số $y=\dfrac{(2m-n)x^2+mx+1}{x^2+mx+n-6}$ ($m, n$ là tham số) nhận trục hoành và trục trung  làm hai đường tiệm cận. Tính $m+n$.
	\shortans{9} 
	\loigiai{ 
		Theo giả thiết ta có  $\heva{&2m-n=0\\&n-6=0}\Leftrightarrow\heva{&m=3\\&n=6.}$\\
	 Vậy $m+n=9$.}
\end{bt}

% nguồn: Phan2_Lop_12_Dai_so_Chuong_1
\begin{bt}%[2D1H4-2]%[ma:b5fe4b99]
	Có bao nhiêu giá trị nguyên dương của  tham số $m$ để đồ thị hàm số $y=\dfrac{x+1}{x^2+4x+m}$ có đúng hai đường tiệm cận.
	\shortans{2} 
	 \loigiai{
		Ta có $\lim\limits_{x\to\pm\infty} y=0$ nên đồ thị hàm số có  $1$ đường tiệm cận ngang $y=0$.\\
	Để đồ thị hàm số có đúng hai đường tiệm cận thì đồ thị phải có đúng $1$ tiệm cận đứng $\Leftrightarrow$ Phương  trình $f(x)=x^2+4x+m=0$ có nghiệm kép hoặc có hai nghiệm phân biệt trong đó có $1$ nghiệm $x=-1\Leftrightarrow\hoac{&\Delta’=0\\&\heva{&\Delta’>0\\&f(-1)=0}}\Leftrightarrow\hoac{&4-m=0\\&\heva{&4-m>0\\&1-4+m=0}}\Leftrightarrow\hoac{&m=4\\&m=3.}$}
\end{bt}

% nguồn: Phan2_Lop_12_Dai_so_Chuong_1
\begin{bt}%[2D1H4-2]%[ma:55badf39]
	Cho hàm số $y=\dfrac{(2m+1)x^2+3}{\sqrt{x^4+1}}, m$ là tham số. Tìm giá trị của $m$ để đường tiệm cận ngang của đồ thị hàm số đi  qua điểm $A(1;-3)$.
	 \shortans{-2}
	\loigiai{ 
		Ta có $\lim\limits_{x\to+\infty} y=\lim\limits_{x\to-\infty} y=2m+1\Rightarrow d\colon y=2m+1$ là tiệm cận ngang của đồ  thị hàm số đã cho.\\
	Do  $A(1;-3)\in d\Leftrightarrow 2m+1=-3\Leftrightarrow m=-2$.}
\end{bt}

% nguồn: Phan2_Lop_12_Dai_so_Chuong_1
\begin{bt}%[2D1V4-2]%[ma:14ddde6b]
	Cho hàm số $f(x)=\dfrac{x^2-5x+4}{x^3-bx^2-a^2x+a^2b}$ có đồ thị $(C)$, với $a$ và $b$ là hai tham số nguyên. Hỏi có tất cả bao nhiêu bộ số $(a;b)$ để có  đúng hai đường tiệm cận (nếu chỉ xét tiệm cận đứng và tiệm cận ngang)?
	 \shortans{7}
	 \loigiai{
		Ta  có $f(x)=\dfrac{x^2-5x+4}{x^3-bx^2-a^2x+a^2b}=\dfrac{(x-1)(x-4)}{(x-a)(x+a)(x-b)}$.\\
		Điều  kiện xác định: $\heva{&x\neq\pm a\\&x\neq b.}$\\
		Đồ thị $(C)$ của hàm số $f(x)$ luôn có một đường tiệm  cận ngang nên để đồ thị $(C)$ của hàm số $f(x)$ luôn có hai tiệm cận thì đồ thị $(C)$ phải có đúng một tiệm cận đứng.\\
		\textbf{Trường hợp 1:} Phương trình $(x-a)(x+a)(x-b)=0$ có nghiệm $x=1$ và $x=4$, ta có các bộ $(a;b)$ thỏa mãn là $(1;4)$, $(-1;4)$,  $(4;1)$ và $(-4;1)$.\\
		\textbf{Trường hợp 2:} Phương trình có nghiệm đơn $x=-a=1$ và nghiệm kép $x=a=b$ ta có bộ $(a;b)$ thỏa mãn là  $(-1;-1)$.\\
		 \textbf{Trường hợp 3:} Phương trình có nghiệm đơn $x=-a=4$ và nghiệm kép $x=a=b$ ta có bộ $(a;b)$ thỏa mãn là $(-4;-4)$.\\
		\textbf{Trường hợp 4:} Phương trình có nghiệm bội ba hay $a=b=0$ ta có bộ $(a;b)$  thỏa mãn là $(0;0)$.\\
	Vậy có tất cả $7$ cặp $(a;b)$ thỏa  mãn.}
\end{bt}

% nguồn: Phan2_Lop_12_Dai_so_Chuong_1
\begin{bt}%[2D1V4-2]%[ma:efed3adf]
	Cho hàm số $y=\dfrac{x-2}{x^2-2mx-m-2}$. Biết  với $m=\dfrac{a}{b}$ ($a,b\in\mathbb{N},\dfrac{a}{b}$ tối giản) thì đồ thị hàm số có đúng $2$ đường tiệm cận. Tính $a+b$.
	 \shortans{7}
	\loigiai{ 
		Ta có $\lim\limits_{x\to\pm\infty}\dfrac{x-2}{x^2-2mx-m-2}=0\Rightarrow y=0$ là đường tiệm cận ngang của đồ thị hàm số.\\ 
		 Đồ thị hàm số có đúng $2$ đường tiệm cận khi và chỉ khi đồ thị hàm số có đúng $1$ đường tiệm cận đứng $\Leftrightarrow x^2-2mx-m-2=0\quad(1)$ có nghiệm kép hoặc có hai nghiệm phân biệt trong đó có một nghiệm bằng $2$.
		 $\Leftrightarrow\hoac{&\Delta=0\\&\heva{&\Delta>0\\&f(2)=0}}\Leftrightarrow\hoac{&4m^2+4m+8=0\\&\heva{&4m^2+4m+8>0\\&2^2-2m\cdot 2-m-2=0}}\Leftrightarrow\hoac{&vn\\&m=\dfrac{2}{5}}\Leftrightarrow m=\dfrac{2}{5}$.\\
	 Vậy $a=2,b=5\Rightarrow a+b=7$.}
\end{bt}

% nguồn: Phan2_Lop_12_Dai_so_Chuong_1
\begin{bt}%[2D1V4-2]%[ma:88afee2a]
	Gọi $(C_m)$ là đồ thị của hàm số $y = mx + \dfrac{1}{x}\left(^*\right)$, $m$ là tham số. Tìm $m$ để hàm số $\left(^*\right)$ có cực trị và khoảng cách từ điểm  cực tiểu của $\left(C_m\right)$ đến đường tiệm cận xiên bằng $\dfrac{1}{\sqrt{2}}$.
	\shortans{1} 
	 \loigiai{
		Tập xác định:  $\mathscr{D} = \mathbb{R}$. Ta có
		$y' =  m - \dfrac{1}{x^2}=\dfrac{mx^2-1}{x}$.\\
		Hàm số (*) có cực trị khi $y' = 0$ có hai nghiệm phân biệt khác $0 \Leftrightarrow \dfrac{1}{m} >0 \Leftrightarrow m > 0$.\\ 
		Lúc đó $y' = 0 \Leftrightarrow x_1 = -\dfrac{1}{\sqrt{m}}$, $x_2 = \dfrac{1}{\sqrt{m}}$. Lập  bảng biến thiên để suy ra\\
		Điểm cực  tiểu của $(C_m)$ là $M\left(\dfrac{1}{\sqrt{m}}; 2\sqrt{m}\right)$.\\
		Tiệm cận xiên $\Delta\colon  y = mx \Leftrightarrow mx - y = 0$.\\
		Theo đề 
		\begin{eqnarray*}
			&&\mathrm{d}\left(M, \Delta\right) = \dfrac{1}{2}\\ &\Leftrightarrow&  \dfrac{\left| \sqrt{m} - 2\sqrt{m} \right|}{\sqrt{m^2 + 1}} = \dfrac{1}{2}\\ &\Leftrightarrow& \dfrac{\sqrt{m}}{\sqrt{m^2 + 1}} = \dfrac{1}{\sqrt{2}}\\
			&\Leftrightarrow& 2m = m^2 +  1\\
			& \Leftrightarrow& m =  1.
		\end{eqnarray*}
		Vậy giá trị $m$ cần tìm  là $m=1$.
	 }
\end{bt}

% nguồn: Phan2_Lop_12_Dai_so_Chuong_1
\begin{bt}%[2D1V4-2]%[ma:1fe9ee82]
	Cho hàm số $y = \dfrac{mx^2 + (3m^2 - 2)x - 2}{x + 3m}~\left(1\right)$ với $m$ là tham số thực. Tổng tất cả các giá trị của $m$ để góc giữa hai đường tiệm cận của  đồ thị hàm số $(1)$ bằng $45^\circ$.
	\shortans{0} 
	 \loigiai{
		Ta  có $y = \dfrac{mx^2 + (3m^2 - 2)x - 2}{x + 3m} = mx - 2 + \dfrac{6m - 2}{x + 3m}$.\\
		Khi  $m = \dfrac{1}{3}$ thì đồ thị hàm số không tồn tại hai tiệm cận.\\
		Khi $m \ne \dfrac{1}{3}$ thì đồ thị hàm số có tiệm cận xiên là  $\Delta_1\colon y = mx - 2 \Leftrightarrow mx - y - 2 = 0$.\\
		Tiệm  cận đứng là $\Delta_2\colon x = -3m \Leftrightarrow x + 3m = 0$.\\
		Góc giữa hai đường tiệm cận của đồ thị hàm số $(1)$ bằng $45^\circ$ nên\\ 
		\[\cos 45^\circ =  \dfrac{\left|\overrightarrow{n}_1\cdot \overrightarrow{n}_2\right|}{\left|\overrightarrow{n}_1\right|\cdot\left|\overrightarrow{n}_2\right|}=\dfrac{\left|m\right|}{\sqrt{m^2+1}}=\dfrac{\sqrt{2}}{2}\Leftrightarrow m=\pm 1.\]
		Vậy giá trị $m$ cần tìm là  $m = \pm 1$.
	} 
\end{bt}

% nguồn: Phan2_Lop_12_Dai_so_Chuong_1
\begin{bt}%[2D1V4-2]%[ma:1c301d29]
	Tính tổng giá trị $a+b$ của  tham số $m=\dfrac{a}{b}$ (tối giản và $a, b$ nguyên dương) sao cho đồ thị hàm số $y = \dfrac{2x^2 + 3mx - m + 2}{x - 1}$ có tiệm cận xiên tạo với các trục tọa độ một tam giác có diện tích bằng $4$ (làm tròn kết quả đến hàng phần trăm).
	 \shortans{5}
	 \loigiai{
		 Ta có
		$y = \dfrac{2x^2 + 3mx - m + 2}{x  - 1} = 2x + 3m + 2 + \dfrac{2m + 4}{x - 1}$.\\
		Với $m = -2$ thì hàm số trở thành $y = 2x - 4$. Đây là hàm số bậc nhất nên  không có tiệm cận.\\
		Với $m \ne  2$ thì
		$\lim\limits_{x \to \pm\infty}  \left[f(x) - \left( 2x + 3m + 2 \right) \right] = \lim\limits_{x \to \pm\infty} \dfrac{2m + 4}{x + 2} = 0$.\\
		Do đó với $m \ne 2$ thì đồ thị hàm số có tiệm cận xiên là $\Delta\colon  y = 2x + 3m + 2$.\\
		$\Delta$ cắt trục hoành  tại điểm $A\left(\dfrac{-3m - 2}{2}; 0\right)$ và cắt trục tung tại $B(0; 3m + 2)$.\\
		Từ đó $OA = \left| \dfrac{-3m - 2}{2} \right| = \left| \dfrac{3m + 2}{2} \right|$, $OB = \left|3m  + 2\right|$.\\
		Theo giả thiết 
		\begin{eqnarray*}
			&&S_{\Delta OAB}  = 4\\
			&\Leftrightarrow&\dfrac{1}{2} OA \cdot  OB = 4\\
			&\Leftrightarrow&\dfrac{|3m + 2|}{2} \cdot  |3m + 2| = 8\\
			&\Leftrightarrow&(3m  + 2)^2 = 16\\
			&\Leftrightarrow&\hoac{&  3m + 2 = -4\\&3m + 2 = 4}\\
			&\Leftrightarrow&\hoac{&m  = -2 \text{ (loại)}\\&m = \dfrac{2}{3}.}
		\end{eqnarray*}
		Vậy giá trị $m$ cần tìm là $m=\dfrac{2}{3} \Rightarrow a+b=5$. 
	} 
\end{bt}

% nguồn: Phan2_Lop_12_Dai_so_Chuong_1
\begin{bt}%[2D1V4-2]%[ma:91c94351]
	Cho hàm số $y=\dfrac{2mx+m}{x-1}$. Có bao  nhiêu giá trị thực của tham số $m$ để đường tiệm cận đứng, tiệm cận ngang của đồ thị hàm số cùng với hai trục tọa độ tạo thành một hình chữ nhật có diện tích bằng $8$.
	\shortans{2} 
	 \loigiai{
		Điều kiện để đồ thị hàm số có tiệm  cận là $-2m-m\neq 0\Leftrightarrow m\neq 0$.\\
		Khi đó, đồ thị  hàm số có\\
		Tiệm cận đứng: $x=1$, song song với $Oy$  và cắt $Ox$ tại điểm $A(1;0)$.\\
		Tệm  cận ngang: $y=2m$, song song với $Ox$ và cắt $Oy$ tại điểm $B(2m;0)$.\\
		Diện tích hình chữ nhật tạo bởi hai đường tiệm cận cùng với hai trục tọa độ là 
	 $$S=OA\cdot OB=1\cdot|2m|=8\Rightarrow m=\pm 4.$$}
\end{bt}

% nguồn: Phan2_Lop_12_Dai_so_Chuong_1
\begin{bt}%[2D1V4-2]%[ma:c9112a00]
	Có bao nhiêu giá trị thực của tham số $m$ để đồ thị hàm số $y=\dfrac{x+1}{x^2+2mx+3m+4}$ có đúng một đường tiệm cận đứng là 
	\shortans{3} 
	\loigiai{ 
		Để đồ thị hàm số $y=\dfrac{x+1}{x^2+2mx+3m+4}$ có đúng một đường tiệm cận đứng thì phương trình $x^2+2mx+3m+4=0$ có duy nhất một nghiệm hoặc phương trình $x^2+2mx+3m+4=0$ có hai nghiệm phân biệt, trong đó có một nghiệm  $x=-1$ và một nghiệm $x$ khác $-1$.\\
		\textbf{Trường hợp 1:}  Phương trình $x^2+2mx+3m+4=0$ có duy nhất một nghiệm.\\
		Suy ra $\Delta’=0\Leftrightarrow m^2-3m-4=0\Leftrightarrow\hoac{&m=-1\\&m=4.}$\\ 
		\textbf{Trường hợp 2:} Phương trình có một nghiệm bằng $-1$, một nghiệm  khác $-1$.\\
		Khi đó ta có $(-1)^2+2\cdot  (-1)m+3m+4=0\Leftrightarrow m=-5$.\\
		Thử lại thấy với $m=-5$ phương trình có $2$ nghiệm phân biệt thỏa mãn điều  kiện.\\
	Vậy  $m\in\{-5;-1; 4\}$ nên có ba giá trị.}
\end{bt}


%%%%%%%%%%%%%%%% TỰ LUẬN %%%%%%%%%%%%%%%%

% nguồn: Phan2_Lop_12_Dai_so_Chuong_1
\begin{bt}%[2D1H4-2]%[ma:2449ffb6]
	Tìm tham số $m$ để đồ  thị hàm số.
	\begin{listEX}[1]
		 \item $y=\dfrac{4x+1}{mx-1}$ không có tiệm cận đứng.
		 \item $y=\dfrac{x^2+x-2}{x^2-2x+m}$ có hai tiệm cận đứng.
		\item $y=\dfrac{3x-1}{x-m}$ có đường tiệm cận đứng là $x=5$. 
		\item  $y=\dfrac{(m+1)x-5m}{2x-m}$ có tiệm cận ngang là $y=1$.
		\item $y=\dfrac{x+3}{x^2-2mx+9}$ có đúng hai đường tiệm cận gồm một  tiệm cận đứng và một tiệm cận ngang.
		\item $y=\dfrac{x-1}{x^2-(2m+1)x+m^2-3}$ có đúng hai đường tiệm cận gồm một tiệm cận đứng  và một tiệm cận ngang.
	\end{listEX}
	 \loigiai{
		\begin{enumerate}
			\item Đồ thị hàm số không có tiệm cận  đứng khi $\hoac{&\heva{&m\neq 0\\&4\cdot (-1)-1\cdot m=0}\\&m=0}\Leftrightarrow\hoac{&m=-4\\&m=0.}$
			\item Để đồ thị hàm số có hai đường tiệm cận đứng thì phương trình $x^2-2x+m=0$ có hai  nghiệm phân biệt khác nghiệm của phương trình $x^2+x-2=0$.\\
			Hay phương trình $f(x)=x^2-2x+m=0$ có hai nghiệm phân biệt  khác $1$ và $-2$.\\
			Khi  đó ta có $\heva{&\Delta’>0\\&f(1)\neq 0\\&f(-2)\neq 0}\Leftrightarrow\heva{&1-m>0\\&m-1\neq 0\\&8+m\neq 0}\Leftrightarrow\heva{&m<1\\&m\neq 1\\&m\neq-8}\Leftrightarrow-8\neq m<1$.\\
			\item Điều kiện để đồ  thị hàm số có tiệm cận đứng là $-3m+1\neq 0\Leftrightarrow m\neq\dfrac{1}{3}$.\\
			Đồ thị hàm số có tiệm cận đứng là $x=m$.\\ 
			Theo yêu  cầu bài toán suy ra $m=5$ (thoả mãn).
			\item Điều kiện để đồ thị hàm số  có tiệm cận ngang là $-m(m+1)+10m\neq 0$.\\
			Tiệm cận  ngang là $y=\dfrac{a}{c}=\dfrac{m+1}{2}$ theo đề bài $\dfrac{m+1}{2}=1\Leftrightarrow m=1$ (thoả mãn).\\
			\item Ta có $\lim\limits_{x\to\pm\infty} y=\lim\limits_{x\to\pm\infty}\dfrac{\dfrac{1}{x}+\dfrac{3}{x^2}}{1-\dfrac{2m}{x}+\dfrac{9}{x^2}}=0$ nên đồ thị hàm số đã cho luôn có một tiệm cận ngang là $y=0$. Do đó đồ thị hàm số đã cho có hai đường  tiệm cận thì đồ thị hàm số đã cho có đúng một tiệm cận đứng.\\
			Phương trình $f(x)=x^2-2mx+9=0$  có nghiệm kép hoặc có hai nghiệm phân biệt trong đó có một nghiệm bằng $-3\Leftrightarrow\hoac{&\Delta’=m^2-9=0\\&\heva{&\Delta’=m^2-9>0\\&f(-3)=0}}\Leftrightarrow\hoac{&m=3\\&m=-3\\&\heva{&m^2-9>0\\&m=-3}}\Leftrightarrow\hoac{&m=3\\&m=-3}$.
			\item Đồ thị hàm số $y=\dfrac{x-1}{x^2-(2m+1)x+m^2-3}$ có $1$ tiệm cận ngang là $y=0$.\\ 
			Đồ thị hàm số $y=\dfrac{x-1}{x^2-(2m+1)x+m^2-3}$ có  đúng hai đường tiệm cận\\
			 $ \Leftrightarrow $ Đồ thị hàm số $y=\dfrac{x-1}{x^2-(2m+1)x+m^2-3}$ có đúng $1$ tiệm cận đứng\\
			$ \Leftrightarrow $ Phương trình $x^2-(2m+1)x+m^2-3=0$ có nghiệm kép hoặc  phương trình $x^2-(2m+1)x+m^2-3=0$ có hai nghiệm phân biệt trong đó có một nghiệm bằng $1$
			 $$\Leftrightarrow\hoac{&\Delta=0\\&\heva{&\Delta>0\\&1-(2m+1)+m^2-3=0}}
			\Leftrightarrow\hoac{&(2m+1)^2-4(m^2-3)=0\\&\heva{&(2m+1)^2-4(m^2-3)>0\\&m^2-2m-3=0}} 
			 \Leftrightarrow\hoac{&m=-\dfrac{13}{4}\\&\heva{&m=3\\&m=-1.}}$$
	\end{enumerate}}
\end{bt}

% nguồn: Phan2_Lop_12_Dai_so_Chuong_1
\begin{bt}%[2D1H4-2]%[ma:05b0025e]
	Xác định tiệm cận đứng và tiệm cận xiên của đồ thị hàm số cho bởi công thức  sau:
	\begin{listEX}[2]
		\item  $y=\dfrac{x-2}{x^2-mx+1}$ có hai tiệm cận đứng.
		\item  $y=\dfrac{2x^2-3x+m}{x-m}$ có tiệm cận xiên.
	\end{listEX}
	\loigiai{ 
		\begin{enumerate}
			\item Đồ thị hàm số có hai đường tiệm cận đứng khi và chỉ khi phương trình $g(x)=x^2-mx+1$ có hai nghiệm phân biệt khác $2\Leftrightarrow\heva{&a=1\neq  0\\&\Delta=m^2-4>0\\&g(2)=2^2-2m+1\neq 0}\Leftrightarrow\heva{&\hoac{&m <-2\\&m>2}\\&m\neq\dfrac{5}{2}.}$
			\item Đồ thị hàm số có đường tiệm  xiên đứng khi và chỉ khi phương trình $g(x)=2x^2-3x+m=0$ không có nghiệm $x=m$ tức là $g(m)\neq 0\Leftrightarrow 2m^2-2m\neq 0\Leftrightarrow\heva{&m\neq 0\\&m\neq 1.}$
	\end{enumerate}}
\end{bt}

% nguồn: Phan2_Lop_12_Dai_so_Chuong_1
\begin{bt}%[2D1V4-2]%[ma:22816356]
	Có bao nhiêu giá trị $m$  nguyên thuộc đoạn $[-2025; 2025]$ để đồ thị hàm số $y = \dfrac{x + \sqrt{x}}{x^2 - 2(m - 1)x + m}$ có hai đường tiệm cận đứng?
	 \loigiai{
		Số đường tiệm cận đứng của đồ thị hàm  số là số nghiệm của hệ $\heva{& x \ge 0 \\ & x + \sqrt{x} \ne 0 \\ & x^2 - 2(m - 1)x + m = 0} \ (*)$\\
		Để đồ thị của hàm số đã cho có hai đường tiệm cận đứng $\Leftrightarrow (*)$ có hai nghiệm phân biệt  $\heva{& x > 0 \\ & x^2 - 2(m - 1)x + m = 0}$ có hai nghiệm phân biệt $\Leftrightarrow x^2-2(m - 1)x + m = 0$ có hai nghiệm phân biệt $x_1$; $x_2$ thỏa mãn $0<x_1<x_2$. \\
		$\Leftrightarrow \heva{& \Delta'=(m - 1)^2 - m > 0 \\ & S>0 \\ & P>0 } \Leftrightarrow \heva{& m^2-3m+1>0  \\ & m-1>0 \\ & m>0} \Leftrightarrow \heva{& \hoac{& m < \dfrac{3-\sqrt{5}}{2} \\ & m > \dfrac{3 + \sqrt {5}}{2}} \\ & m > 1} \Leftrightarrow m>\dfrac{3+\sqrt {5}}{2}$. \\
		Mặt khác $-2025 \le m  \le 2025$ và $m \in \mathbb{Z}$.\\
		Do đó $m  \in \{3; 4;\cdots;2025 \}$.\\
		Vậy có $2023$ số nguyên  $m$ phải tìm.
	 }
\end{bt}
